\subsection{DERIVATIONS OF FREE LIE ALGEBRAS}

PROPOSITION 8. Let X be a set, let M be an L(X)-module and let d be a mapping of X into M. There exists one and only one linear mapping D of L(X) into M extending d and satisfying the relation:

\[
\mathrm{D} ([ a, a ^ {\prime} ]) = a. \mathrm{D} (a ^ {\prime}) - a ^ {\prime}. \mathrm{D} (a) \quad for a, a ^ {\prime} in \mathrm{L} (\mathrm{X}).\tag{20}
\]

We define a Lie algebra g with underlying module M × L(X) by means of the bracket

\[
[ (m, a), (m ^ {\prime}, a ^ {\prime}) ] = (a. m ^ {\prime} - a ^ {\prime}. m, [ a, a ^ {\prime} ]),\tag{21}
\]

for \(a, a'\) in \(\mathrm{L}(\mathrm{X})\) and \(m, m'\) in \(\mathbf{M}\) (Chapter I, § 1, no. 8). Let \(f\) be the homomorphism of \(\mathrm{L}(\mathrm{X})\) into \(\mathfrak{g}\) such that \(f(x) = (d(x), x)\) for all \(x\) in \(\mathbf{X}\) ; let \(f(a) = (\mathrm{D}(a), u(a))\) for all \(a\) in \(\mathrm{L}(\mathrm{X})\) . By formula (21), \(u\) is a homomorphism of \(\mathrm{L}(\mathrm{X})\) into itself; as \(u(x) = x\) for \(x\) in \(\mathbf{X}\) , \(u(a) = a\) for all \(a\) in \(\mathrm{L}(\mathrm{X})\) , whence

\[
f (a) = (\mathrm{D} (a), a).\tag{22}
\]

By (21) and (22), relation (20) then follows from \(f([a, a']) = [f(a), f(a')]\) .

Conversely, let \(D'\) be a mapping of \(\mathrm{L}(\mathrm{X})\) into M satisfying relation \((20')\) analogous to (20) and extending d.~Let \(f'(a) = (\mathrm{D}'(a), a)\) for \(a \in \mathrm{L}(\mathrm{X})\) ; by \((20')\) and \((21)\) , \(f'\) is a homomorphism of \(\mathrm{L}(\mathrm{X})\) into g, coinciding with f on X, whence \(f' = f\) and \(D' = D\) .

COROLLARY. Every mapping of \(\mathbf{X}\) into \(\mathrm{L}(\mathbf{X})\) can be extended uniquely to a derivation of \(\mathrm{L}(\mathbf{X})\) .

When M is equal to L(X) with the adjoint representation, relation (20) means that D is a derivation.

